\documentclass{article}
\usepackage[top=1.5cm, bottom=1.5cm, left=2cm, right=1.5cm]{geometry}
\usepackage[hidelinks,unicode]{hyperref}
\usepackage[dethi]{ex_test} 
\usepackage{environ} 
\usepackage{lastpage} 
\usepackage{fancyhdr} 
\usepackage{tcolorbox}
\tcbuselibrary{skins,xparse,breakable,raster,listings,listingsutf8,documentation,many}
\usepackage{xcolor}
\usepackage{amsmath}
\usepackage{graphicx}

\def\vec{\overrightarrow}
\newcommand{\hoac}[1]{\left[\begin{aligned}#1\end{aligned}\right.}
\newcommand{\heva}[1]{\left\{\begin{aligned}#1\end{aligned}\right.}
\newcommand{\hetde}{\centerline{\rule[0.5ex]{2cm}{1pt} HẾT \rule[0.5ex]{2cm}{1pt}}}

\newcommand{\tentruong}{THPT Nguyễn Hữu Cảnh}
\newcommand{\tengv}{GV Nguyễn Văn Sang}
\newcommand{\tenkythi}{Ôn Kiểm Tra GHKII 2025}
\newcommand{\tenmonthi}{Môn: Toán - Khối 12}
\newcommand{\thoigian}{90}
\newcommand{\made}{101}

\newcommand{\tieude}[3]{
	\noindent
	% Trái
	\begin{minipage}[b]{7cm}
		\centerline{\textbf{\fontsize{13}{0}\selectfont \tentruong}}
		\centerline{\textbf{\fontsize{13}{0}\selectfont \tengv}}
		\centerline{(\textit{Đề thi có #1\ trang})}
		\centerline{\textit{ }}
	\end{minipage}\hspace{1.5cm}
	% Phải
	\begin{minipage}[b]{9cm}
		\centerline{\textbf{\fontsize{13}{0}\selectfont \tenkythi}}
		\centerline{\textbf{\fontsize{13}{0}\selectfont \tenmonthi}}
		\centerline{\textit{\fontsize{12}{0}\selectfont Thời gian làm bài \thoigian \ phút }}
		% Đã sửa lỗi Overfull \hbox bằng cách ngắt dòng
		\centerline{\textit{\fontsize{12}{0}\selectfont (Đề có 11 câu trắc nghiệm, 5 câu đúng sai, 6 câu điền đáp án)}}
	\end{minipage}
\vspace{0.5em}

	\begin{minipage}[b]{13cm}
		\vspace*{.75cm}
		\textbf{Họ và tên thí sinh: }{\tiny\dotfill}\\
		\textbf{Số Báo Danh: }{\tiny\dotfill}
	\end{minipage}
	\begin{minipage}[b]{8cm}
		\hspace*{1cm}\fbox{\bf Mã đề thi #3}
	\end{minipage}
}
% Định nghĩa lại chantrang để thiết lập footer cho fancyhdr
\newcommand{\chantrang}[2]{
    \fancyhf{} % Xóa tất cả các header và footer mặc định
    \rfoot{Trang \thepage/#1 $-$ Mã đề #2} % Đặt nội dung vào chân trang bên phải
    \lfoot{} % Để trống chân trang bên trái
    \cfoot{} % Để trống chân trang ở giữa
}
\renewcommand{\headrulewidth}{0pt}
% --- Các định nghĩa ---
\newcommand{\caulc}{
	\begin{flushleft}
{\color{red}\sffamily\bfseries Phần I. Câu hỏi trắc nghiệm nhiều phương án lựa chọn}\\
\textit{Thí sinh trả lời từ câu 1 đến câu 12. Mỗi câu hỏi thí sinh chỉ chọn một phương án.}
	\end{flushleft}
\setcounter{ex}{0}


	\setcounter{ex}{0}
}
\newcommand{\cauds}{
	\begin{flushleft}
{\color{red}\sffamily\bfseries Phần II. Câu hỏi trắc nghiệm đúng sai}\\
\textit{Thí sinh trả lời từ câu 1 đến câu 4. Trong mỗi ý a), b), c), d) ở mỗi câu thí sinh chọn đúng hoặc sai.} 
	\end{flushleft}
	\setcounter{ex}{0}
}
\newcommand{\caukq}{
	\begin{flushleft}
{\color{red}\sffamily\bfseries Phần III. Câu hỏi trắc nghiệm trả lời ngắn}\\
\textit{Thí sinh trả lời từ câu 1 đến câu 6.} 	\end{flushleft}
	\setcounter{ex}{0}
} 

% === CƠ CHẾ ẨN/HIỆN LỜI GIẢI ===
\newif\ifhienloigiai

\newcommand{\hienloigiai}{
    \hienloigiaitrue
    \makeatletter
    \def\colorEX{\color{red}}
    \renewcommand{\TrueEX}{\stepcounter{dapan}\squareEX{\bf\Alph{dapan}}\ignorespaces}
    \makeatother
    \typeout{>>> CHE DO DAP AN DA DUOC KICH HOAT.}
}

\newcommand{\anloigiai}{
    \hienloigiaifalse
    \makeatletter
    \def\colorEX{}
    \renewcommand{\TrueEX}{\FalseEX}
    \makeatother
    \typeout{>>> CHE DO DE THI DA DUOC KICH HOAT.}
}

\RenewEnviron{loigiai}{%
    \ifhienloigiai
        \par\noindent\textbf{\loigiaiEX}
        \BODY
    \fi
    \makeatletter
    \xdef\keyEX{}%
    \setcounter{numTrueSol}{0}%
    \makeatother
}

\anloigiai
\usepackage{polyglossia}
\setmainlanguage{vietnamese}
\usepackage{fontspec}
\setmainfont{Times New Roman}

% Load Premium Design Package
\usepackage{premium_design}

\begin{document}
\pagestyle{fancy} 
\tieude{\pageref{LastPage}}{0}{\made} 
\chantrang{\pageref{LastPage}}{\made}

\section{Chương 1: Hàm Số Lũy Thừa (Premium Design)}

\begin{dn}{Hàm số lũy thừa}{dn:luythua}
    Hàm số $y=x^\alpha$ với $\alpha \in \mathbb{R}$ được gọi là hàm số lũy thừa.
\end{dn}

\begin{dl}{Đạo hàm}{dl:daoham}
    $(x^\alpha)' = \alpha x^{\alpha-1}$.
\end{dl}

\begin{vd}{Tìm tập xác định}{}
    Tìm tập xác định của hàm số $y=(x-1)^{\frac{1}{3}}$.
\end{vd}

\begin{bt}{Tính đạo hàm (Có trắc nghiệm)}{}
    Tính đạo hàm của $y=x^{\pi}$.
    \choiceTF
	{$\dfrac{9}{\sqrt{35}}$}
	{$-\dfrac{9}{\sqrt{35}}$}
	{$-\dfrac{9}{2\sqrt{35}}$}
	{\True $\dfrac{9}{2\sqrt{35}}$}
	\loigiai{
		Ta có $\overrightarrow{AB} = (1;5;-2)$, $\overrightarrow{AC} = (5;4;-1)$.\\
		Do đó
		$$ \cos \widehat{BAC} = \dfrac{\overrightarrow{AB}\cdot\overrightarrow{AC}}{\left|\overrightarrow{AB} \right|\cdot \left| \overrightarrow{AC}\right|} = \dfrac{5 + 20 + 2}{\sqrt{30}\cdot\sqrt{42}} = \dfrac{9}{2\sqrt{35}}. $$
	}
\end{bt}

\section{Đề Thi Mẫu (Standard ex)}

\subsection{Phần I: Trắc nghiệm nhiều lựa chọn}
\caulc
\Opensolutionfile{ans}[ans-phanI]

\begin{ex}
Phương trình tiệm cận đứng của đồ thị hàm số $y=\dfrac{x+1}{x+2}$ là:
\choice
    {$x=-1$}
    {\True $x=-2$}
    {$x=2$}
    {$y=1$}
\loigiai{
Ta có $\displaystyle \lim_{x\to -2^+}\frac{x+1}{x+2} = -\infty$, nên $x=-2$ là tiệm cận đứng.
}
\end{ex}
\Closesolutionfile{ans}

\subsection{Phần II: Trắc nghiệm đúng sai}
\cauds
\Opensolutionfile{ans}[ans-phanII]

\begin{ex}
Phương trình tiệm cận đứng của đồ thị hàm số $y=\dfrac{x+1}{x+2}$ là:
\choiceTFt
    {$x=-1$}
    {\True $x=-2$}
    {$x=2$}
    {$y=1$}
\loigiai{
Ta có $\displaystyle \lim_{x\to -2^+}\frac{x+1}{x+2} = -\infty$, nên $x=-2$ là tiệm cận đứng.
}
\end{ex}
\Closesolutionfile{ans}

\subsection{Phần III: Trả lời ngắn}
\caukq
\Opensolutionfile{ans}[ans-phanIII]

\begin{ex}
Biết hàm số $f(x)=\dfrac{4x^2 + 6x + 4}{(x + 7)(x + 10)}$ có nguyên hàm $F(x)=mx + n\ln|x + 7| + p\ln|x + 10| + C$. Tính giá trị $m + n + p$, làm tròn đến phần nguyên.
\shortans[oly]{$-58$}
\loigiai{
    Ta có:
    \[
    \int \dfrac{4x^2 + 6x + 4}{(x + 7)(x + 10)} \, dx = 4 x + \frac{158}{3} \ln|x + 7| - \frac{344}{3} \ln|x + 10| + C
    \]
    Do đó $m=4, n = \frac{158}{3}, p = -\frac{344}{3}$.
    Tổng $m + n + p = 4 + \frac{158}{3} - \frac{344}{3} = 4 - \frac{186}{3} = 4-62=-58$.
}
\end{ex}
\Closesolutionfile{ans}

\section*{BẢNG ĐÁP ÁN}
\noindent\textbf{A. ĐÁP ÁN PHẦN I}
\inputansbox{10}{ans-phanI}

\noindent\textbf{B. ĐÁP ÁN PHẦN II}
\inputansbox{2}{ans-phanII}

\noindent\textbf{C. ĐÁP ÁN PHẦN III}
\inputansbox{2}{ans-phanIII}

\end{document}
